\section*{第 \ref{ch:g12:light-as-wave} 章——光的波动性：衍射与干涉}

\begin{solution}{exo:g12:light-as-wave:1}
$f = c/\lambda$：红光 $\num{3.00e8}/\num{7.0e-7} \approx
\qty{4.3e14}{Hz}$；紫光 $\approx \qty{7.5e14}{Hz}$——紫光
\omterm{def:g10:signals-and-waves:frequency}{频率}更高（\omterm{def:g12:mechanical-waves:wavelength}{波长}更短，速率同为 $c$）。
\end{solution}

\begin{solution}{exo:g12:light-as-wave:2}
$\theta = \num{5.32e-7}/\num{2.0e-4} = \qty{2.7e-3}{rad}$;
$L = 2\theta D = 2 \times \num{2.66e-3} \times 1.5 \approx
\qty{8.0}{mm}$.
\end{solution}

\begin{solution}{exo:g12:light-as-wave:3}
\omterm{def:g10:signals-and-waves:sound}{声音}：$\lambda = 340/425 = \qty{0.80}{m}$，则 $\lambda/a = 1.0$——
\omterm{def:g12:light-as-wave:diffraction}{衍射}最强，\omterm{def:g10:signals-and-waves:sound}{声音}充满走廊。光：
$\lambda/a \approx \num{5.5e-7}/0.80 \approx \num{7e-7}$ ---
扩散可忽略不计：光沿直线传播，拐角即可将其挡住。
\end{solution}

\begin{solution}{exo:g12:light-as-wave:4}
在 $M$ 处：$\delta = \qty{6.0}{cm} = 3\lambda$，\omterm{prop:g12:light-as-wave:conditions}{相长干涉}——明亮。
在 $N$ 处：$\delta = \qty{5.0}{cm} = (2 + \tfrac12)\lambda$，相消
——（近乎）暗。
\end{solution}

\begin{solution}{exo:g12:light-as-wave:5}
$i = \lambda D/a = \num{6.33e-7} \times 1.8/\num{2.5e-4} \approx
\qty{4.6}{mm}$。在 $\pm\qty{1.0}{cm}$ 之内：$|k| \leq 10/4.6 = 2.2$，
则 $k = -2, \dots, 2$：共五条亮条纹。
\end{solution}

\begin{solution}{exo:g12:light-as-wave:6}
$d = 2\lambda D/L = 2 \times \num{6.5e-7} \times 1.6/0.028 \approx
\qty{74}{\micro m}$.
\end{solution}

\begin{solution}{exo:g12:light-as-wave:7}
$i = \qty{4.7}{mm}$，则 $\lambda = i\,a/D = \num{4.7e-3} \times
\num{2.0e-4}/1.5 \approx \qty{6.3e-7}{m} = \qty{630}{nm}$：红光。
测量十个\omterm{def:g12:light-as-wave:fringes}{条纹间距}，可使刻度尺的读数误差降为原来的十分之一。
\end{solution}

\begin{solution}{exo:g12:light-as-wave:8}
两盏灯的\omterm{def:g10:signals-and-waves:frequency}{频率}（大致）相同，但\emph{不
同步}：每盏灯发出一段段短促的波列，相位随机跳变，因此
两光源不相干，瞬时条纹
每秒移动数十亿次——人眼看到的只是平均效果：均匀一片。
杨用\emph{同一}列波照亮两条缝：两束光
每次相位跳变都一同继承，故始终保持同步。
\end{solution}

\begin{solution}{exo:g12:light-as-wave:9}
第一条至第七条亮纹之间的距离 $= 6i = \qty{18.0}{mm}$，则 $i = \qty{3.0}{mm}$；
$\lambda = i\,a/D = \num{3.0e-3} \times \num{3.0e-4}/1.4 \approx
\qty{6.4e-7}{m} \approx \qty{640}{nm}$.
\end{solution}

\begin{solution}{exo:g12:light-as-wave:10}
$\sin\theta = k\lambda/d = 0.396\,k$: $k = 1$, $\theta = 23^\circ$;
$k = 2$，$\sin\theta = 0.79$，$\theta = 52^\circ$；$k = 3$ 时需要
$\sin\theta = 1.19 > 1$：不可能。最高级次 $k = 2$。
\end{solution}

\begin{solution}{exo:g12:light-as-wave:11}
$\theta = \num{5.5e-7}/0.10 = \qty{5.5e-6}{rad}$。两盏车灯
所张角为 $1.5/\num{1.0e5} = \qty{1.5e-5}{rad}$，约为
\omterm{def:g12:light-as-wave:diffraction}{衍射}模糊角的三倍：能分辨，刚好。
\end{solution}

\begin{solution}{exo:g12:light-as-wave:12}
前、后表面的反射光发生\omterm{prop:g12:light-as-wave:conditions}{干涉}；\omterm{def:g10:light-spectra:white}{白光}下，每种厚度
使部分\omterm{def:g12:mechanical-waves:wavelength}{波长}相长、部分\omterm{def:g12:mechanical-waves:wavelength}{波长}相消，因而呈现一种
颜色。薄膜厚度只与高度有关，因此等色点
构成水平条带。排液使薄膜变薄，因此
产生某一颜色的薄膜厚度位置越来越低：
条带缓慢下移。
\end{solution}

\begin{solution}{exo:g12:light-as-wave:13}
$i = \lambda D/a$：紫光 $\num{4.0e-7} \times 1.5/\num{2.0e-4} =
\qty{3.0}{mm}$；红光 $\qty{5.6}{mm}$。中央：各色均亮——
形成一条白色条纹。离中央数毫米处，各色条纹
已彼此错开：呈现\omterm{def:g12:light-as-wave:iridescence}{虹彩}条纹。
更远处，各色亮纹处处重叠，呈均匀的白色模糊。
\end{solution}

\begin{solution}{exo:g12:light-as-wave:14}
$\theta = \num{5.5e-7}/2.4 = \qty{2.3e-7}{rad}$；在 \qty{250}{km}
处，即 $\num{2.3e-7} \times \num{2.5e5} \approx \qty{6}{cm}$。它能
数清汽车、发现行人，但报纸的字（毫米大小的字母）
比\omterm{def:g12:light-as-wave:diffraction}{衍射}极限小二十倍：此说法纯属谣传。
\end{solution}

\begin{solution}{exo:g12:light-as-wave:15}
\omterm{def:g10:signals-and-waves:frequency}{频率}由光源决定，保持不变。速率降为
$c/n$，故 $\lambda' = \lambda/n$，且
$i' = i/n = 4.0/1.33 \approx \qty{3.0}{mm}$：条纹变密。
\end{solution}

\begin{solution}{pb:g12:light-as-wave:1}
\textbf{1.} $\theta = \lambda/a = \num{6.33e-7}/\num{1.00e-4} =
\qty{6.3e-3}{rad}$.

\textbf{2.} 半宽 $D\tan\theta \approx D\theta = \lambda D/a$，
则 $L = 2\lambda D/a = 2 \times \num{6.33e-7} \times
3.00/\num{1.00e-4} = \qty{3.8}{cm}$.

\textbf{3.} $(3.8 - 3.7)/3.8 \approx 3\%$：在刻度尺的
精度之内——已验证。

\textbf{4.} $L \propto 1/a$：$a$ 减半则 $L$ 加倍，故
$L = \qty{7.6}{cm}$。缝或障碍物越窄，图样越宽。

\textbf{5.} $\tan(\num{6.33e-3}) = \num{6.33008e-3}$：二者在
约 $\num{1e-5}$ 的相对误差范围内相符——此近似远优于
本题中的任何测量精度。

\textbf{6.} 由\cref{prop:g12:light-as-wave:hair}可知：宽度为 $d$ 的障碍物的\omterm{def:g12:light-as-wave:diffraction}{衍射}
与宽度为 $d$ 的狭缝相同，故 $L = 2\lambda D/d$ 仍然
成立。

\textbf{7.} $d = 2\lambda D/L = 2 \times \num{6.33e-7} \times
3.00/0.054 \approx \qty{70}{\micro m}$.

\textbf{8.} $L = \qty{5.2}{cm} \to d = \qty{73}{\micro m}$;
$L = \qty{5.6}{cm} \to d = \qty{68}{\micro m}$：故
$d = \qty{70 \pm 3}{\micro m}$.

\textbf{9.} $d = \num{3.80e-6}/0.076 = \qty{50}{\micro m}$：友人的
头发更细——且 $L \propto 1/d$，故更细的头发所\omterm{def:g12:light-as-wave:diffraction}{衍射}的图样
更宽。

\textbf{10.} 不能：灯发出的是\omterm{def:g10:light-spectra:white}{白光}（每种\omterm{def:g12:mechanical-waves:wavelength}{波长}形成不同的
图样，且相互模糊），而且是非相干的宽光源（没有单一、清晰的光束可供\omterm{def:g12:light-as-wave:diffraction}{衍射}）。
激光是单色的、定向的。

\textbf{11.} 只有同一列波的复制品才是相干的；两个独立的
光源会失去同步，条纹随之消失
（见\cref{rem:g12:light-as-wave:lamps}）。

\textbf{12.} 中心处 $\delta = 0 = 0 \times \lambda$ 对任意
$\lambda$ 均成立：为\omterm{prop:g12:light-as-wave:conditions}{相长干涉}，与\omterm{def:g12:mechanical-waves:wavelength}{波长}无关。

\textbf{13.} $i = 4.26/10 = \qty{4.3}{mm}$；十个\omterm{def:g12:light-as-wave:fringes}{条纹间距}的总长读数
其刻度尺误差与单个\omterm{def:g12:light-as-wave:fringes}{条纹间距}相同，故 $i$ 的误差
为原来的十分之一。

\textbf{14.} $\lambda = i\,a'/D = \num{4.26e-3} \times
\num{2.50e-4}/2.00 = \num{5.33e-7} \approx \qty{533}{nm}$.

\textbf{15.} $\qty{533}{nm}$ 位于 $\qty{532 \pm 10}{nm}$ 之内：
二者相符；该\omterm{def:g12:mechanical-waves:wavelength}{波长}为绿色。

\textbf{16.} $\tan\theta = 0.43/1.00$，故 $\theta = 23^\circ$。此处
$\tan\theta = 0.43$，而 $\sin\theta = 0.40$：在 $23^\circ$ 处，
以小角近似把 $\sin$、$\tan$ 与 $\theta$ 视为相等，误差
接近 $10\%$——应使用精确的三角函数。

\textbf{17.} $d = \lambda/\sin\theta = \num{6.33e-7}/0.396 \approx
\qty{1.6}{\micro m}$.

\textbf{18.} $\tan\theta = 1.65$, $\theta = 59^\circ$,
$\sin\theta = 0.855$: $d = \num{6.33e-7}/0.855 \approx
\qty{0.74}{\micro m}$.

\textbf{19.} $1.6/0.74 \approx 2.2$：DVD 的\omterm{def:g10:universal-gravitation:satellite}{轨道}密度是其两倍，
其凹坑相应地更小（用\omterm{def:g12:mechanical-waves:wavelength}{波长}更短的
激光读取）——同一张
\qty{12}{cm} 光盘上的数据量多达数倍。

\textbf{20.} 头发 \qty{70 \pm 3}{\micro m}；友人的头发
\qty{50}{\micro m}；绿色激光 \qty{533}{nm}；CD \omterm{def:g10:universal-gravitation:satellite}{轨道}
\qty{1.6}{\micro m}；DVD \omterm{def:g10:universal-gravitation:satellite}{轨道} \qty{0.74}{\micro m}——已知
\omterm{def:g12:mechanical-waves:wavelength}{波长}可将每种图样的几何形状换算为长度，因此\omterm{def:g12:mechanical-waves:wavelength}{波长}为
半微米的光，就是一把刻度为半微米的标尺。
\end{solution}
